\section{The decomposition theorem}\label{app:D.28}

A response summed over sectors of final states equals the full response only if the sector projectors commute with the Hamiltonian, and none of the projectors of \cref{sec:8.6} does so exactly. The notation is that of \cref{sec:8.6}; the adiabatic states, surfaces and couplings are those of \cref{sec:5.1,sec:6.1}.

\subsection{Sectors and the bin bound}

Let $\mathcal{P}_A$ be orthogonal projectors with $\sum_A\mathcal{P}_A = 1$, write $H = H_{\mathrm{d}} + H_{\mathrm{off}}$ with $H_{\mathrm{d}} = \sum_A\mathcal{P}_AH\mathcal{P}_A$, and fix one initial state $|i\rangle$ with $\psi_b = O_b|i\rangle$. A measured spectrum is never resolved better than some bin, and we describe the bin by a shape $w_{\mathrm{b}}(x) \ge 0$ with $\int\dd x\,w_{\mathrm{b}} = 1$, its width $\Delta\omega_{\mathrm{bin}}$, the kernel $f_\Delta(x) = w_{\mathrm{b}}(x/\Delta\omega_{\mathrm{bin}})/\Delta\omega_{\mathrm{bin}}$ and the transform $\hat w_{\mathrm{b}}(s) = \int\dd x\,w_{\mathrm{b}}(x)\,\ee^{-\ii sx}$. Without the factor $2\pi/\Vol$, the binned response of one initial state is $S^{\mathrm{bin}}_{ab}(\omega) = \langle i|O_a^\dagger f_\Delta(\omega + E_i - H)O_b|i\rangle$. Since $f_\Delta(\omega + E_i - H_{\mathrm{d}})$ commutes with every $\mathcal{P}_A$, it has no blocks between different sectors, and
\begin{gather*}
  S^{\mathrm{bin}}_{ab} = \sum_A\bigl\langle i\big|O_a^\dagger\,\mathcal{P}_A\,f_\Delta(\omega + E_i - H_{\mathrm{d}})\,\mathcal{P}_A\,O_b\big|i\bigr\rangle + C^{\mathrm{cross}}_{ab} , \\
  C^{\mathrm{cross}}_{ab} = \bigl\langle O_ai\big|\bigl[f_\Delta(\omega + E_i - H) - f_\Delta(\omega + E_i - H_{\mathrm{d}})\bigr]\big|\psi_b\bigr\rangle .
\end{gather*}
Each sector term is a restriction by orthogonal projectors, positive for $a = b$ and positive semidefinite as a matrix. To bound $C^{\mathrm{cross}}_{ab}$, write the kernel as a time integral, $f_\Delta(\omega + E_i - H) = \int(\dd t/2\pi)\,\hat w_{\mathrm{b}}(\Delta\omega_{\mathrm{bin}}t)\,\ee^{\ii(\omega + E_i)t}\ee^{-\ii Ht}$, and compare the two evolutions with the Duhamel formula, which follows by differentiating $\ee^{\ii Ht}\ee^{-\ii H_{\mathrm{d}}t}$:
\begin{equation*}
  \bigl(\ee^{-\ii Ht} - \ee^{-\ii H_{\mathrm{d}}t}\bigr)\psi_b = -\ii\int_0^t\dd t'\;\ee^{-\ii H(t - t')}\,H_{\mathrm{off}}\,\ee^{-\ii H_{\mathrm{d}}t'}\,\psi_b .
\end{equation*}
With $\mathfrak{h}_{\mathrm{off}} = \sup_{t'}\|H_{\mathrm{off}}\ee^{-\ii H_{\mathrm{d}}t'}\psi_b\|/\|\psi_b\|$ the norm of the right side is at most $|t|\,\mathfrak{h}_{\mathrm{off}}\|\psi_b\|$. Inserting it in the time integral and substituting $y = \Delta\omega_{\mathrm{bin}}t$, the Cauchy--Schwarz inequality gives
\begin{equation}
  |C^{\mathrm{cross}}_{ab}| \le \mathfrak{c}_{\mathrm{b}}\,\frac{\mathfrak{h}_{\mathrm{off}}}{\Delta\omega_{\mathrm{bin}}^2}\,\bigl\|O_a|i\rangle\bigr\|\,\bigl\|O_b|i\rangle\bigr\| , \qquad
  \mathfrak{c}_{\mathrm{b}} = \frac{1}{2\pi}\int\dd y\,|y|\,|\hat w_{\mathrm{b}}(y)| .
  \label{eq:D.28-lemma}
\end{equation}
For a Gaussian bin of standard deviation $\Delta\omega_{\mathrm{bin}}$, $w_{\mathrm{b}}(x) = \ee^{-x^2/2}/\sqrt{2\pi}$, the transform is $\ee^{-y^2/2}$, the integral $\int\dd y\,|y|\,\ee^{-y^2/2} = 2$, and $\mathfrak{c}_{\mathrm{b}} = 1/\pi$, the constant of \cref{eq:8.6-bin}. A rectangular bin has a transform that falls only as $1/y$, and $\mathfrak{c}_{\mathrm{b}}$ diverges: the bound requires bins with smooth edges. Only the state $\psi_b$ that the probe creates enters $\mathfrak{h}_{\mathrm{off}}$, not the norm of $H_{\mathrm{off}}$ over the whole space, which is proportional to the volume for a local coupling.

\subsection{The electron-following projector}

The hypotheses of \cref{sec:8.6} are those under which the space-adiabatic theorem applies. The electronic Hamiltonian $H_{\mathrm{el}}(u)$ must be a smooth function of the nuclear positions; with nuclei of finite size, rigid charge distributions $Z_\kappa F^{\mathrm{ch}}_\kappa$ with $F^{\mathrm{ch}}_\kappa$ the nuclear charge form factor, the electron--nucleus potential is smooth and Spohn and Teufel's regularity condition holds, whereas for point nuclei it fails~\cite[Sect.~2]{spohn2001adiabatic}. The electron--electron repulsion does not depend on $u$ and does not enter. The ground-state energy $E_0(u)$ must be nondegenerate and separated from the rest of the spectrum of $H_{\mathrm{el}}(u)$ by a gap, here at least $\Eg$, on the configurations that the vibrational state visits, up to exponentially small tails. The lattice must be harmonic with perturbative anharmonicity, and $T \ll \Eg$.

In units where the electronic scales are one and the nuclear momenta are measured in units of $(M/m_e)^{1/2}$, the nuclear kinetic energy is $-\kappa_{\mathrm{BO}}^4\nabla^2/2$, the small parameter of the semiclassical expansion is $(m_e/M)^{1/2} = \kappa_{\mathrm{BO}}^2$, and the full Hamiltonian is the quantization of the operator-valued symbol $\tfrac12p^2 + H_{\mathrm{el}}(u)$~\cite[Eq.~(38)]{panati2003space}. Theorem~2.1 of Panati, Spohn and Teufel then provides an orthogonal projector $\Pi$ whose symbol is a series $\sum_j\kappa_{\mathrm{BO}}^{2j}\pi_j$ with leading term the spectral projector of the electronic ground state, $\pi_0(u) = |\Psi_0(u)\rangle\langle\Psi_0(u)|$, and with $[H, \Pi]$ smaller than any power of $\kappa_{\mathrm{BO}}$~\cite[Thm.~2.1]{panati2003space}. The Born--Oppenheimer Hamiltonian grows quadratically in the nuclear momentum while the gap does not, and the theorem is applied below an arbitrary fixed energy through an auxiliary symbol equal to the true one on the region of phase space where the classical energy $\frac12p^2 + E_0(u)$ lies below it~\cite[Sect.~4.5]{panati2003space}. Truncating the series after $j = n$ defines $\mathcal{P}_{\mathrm{ad}}$, and since $(1 - \mathcal{P}_{\mathrm{ad}})H\mathcal{P}_{\mathrm{ad}} = (1 - \mathcal{P}_{\mathrm{ad}})[H, \mathcal{P}_{\mathrm{ad}}]\mathcal{P}_{\mathrm{ad}}$, the leak is of the order of the first omitted term, $\kappa_{\mathrm{BO}}^{2(n+1)}$, in operator norm on the states below that energy. This is \cref{eq:8.6-leak}. Low-lying vibrational states lie below any such energy above the zero-point energy, so the bound applies to them with the same order. The constants depend on the number of slow coordinates, which is why the proof covers a finite crystal.

The first correction to the state follows more simply from the coupled equations \cref{eq:5.1-coupled}. Let a state on the ground surface have the nuclear wavefunction $\Psi^{\mathrm{nuc}}_0$ and the energy $E$, and let $\Psi^{\mathrm{nuc}}_j$ be its component on the excited adiabatic state $\Psi_j$, as in \cref{sec:5.1}. To first order in the coupling, $[T_N + E_j(u) - E]\Psi^{\mathrm{nuc}}_j = -\mathcal{N}_{j0}\Psi^{\mathrm{nuc}}_0$, with $T_N$ the kinetic energy of the nuclei. The derivative coupling dominates $\mathcal{N}_{j0}\Psi^{\mathrm{nuc}}_0$, which is $-M_\kappa^{-1}\vect{d}_{j0}\cdot\nabla\Psi^{\mathrm{nuc}}_0$ by \cref{eq:5.1-N,eq:6.1-d}, and the operator on the left is $E_j(u) - E_0(u) \equiv \Delta E_j$ up to vibrational energies. Hence
\begin{equation}
  \Psi^{\mathrm{nuc}}_j = \frac{\vect{d}_{j0}\cdot\nabla\Psi^{\mathrm{nuc}}_0}{M_\kappa\,\Delta E_j} = \frac{\ii\,\vect{d}_{j0}\cdot\vect{p}\,\Psi^{\mathrm{nuc}}_0}{M_\kappa\,\Delta E_j} ,
  \label{eq:D.28-tilt}
\end{equation}
with $\vect{p} = -\ii\nabla$ the nuclear momentum, a sum over nuclei being understood. It is the first-order term of the unitary of Panati, Spohn and Teufel, written in physical units~\cite[below Eq.~(39)]{panati2003space}.

The same order of the expansion changes the mass. The second-order effective Hamiltonian on the ground surface is~\cite[Eq.~(40)]{panati2003space}, for a real nondegenerate state with vanishing Berry connection,
\begin{equation*}
  h_{\mathrm{BO}} = \frac{p^2}{2M} + E_0(u) + \frac{1}{2M}\sum_{j\neq0}|\vect{d}_{j0}|^2 - \sum_{j\neq0}\frac{(\vect{p}\cdot\vect{d}_{0j})(\vect{d}_{j0}\cdot\vect{p})}{M^2\,\Delta E_j} ,
\end{equation*}
in physical units for one nucleus. The third term is the diagonal correction of Born and Huang (\cref{app:D.13}). The last term is quadratic in the momentum. With $\vect{d}_{j0} = \vect{F}_{j0}/(E_0 - E_j)$ from \cref{eq:5.1-dcoup} it equals $-\vect{p}\cdot\delta M\cdot\vect{p}/2M^2$ with
\begin{equation}
  \delta M = 2\sum_{j\neq0}\frac{\Real\bigl[\vect{F}_{0j}\vect{F}_{j0}\bigr]}{\Delta E_j^3} ,
  \label{eq:D.28-dM}
\end{equation}
a positive tensor, and together with $p^2/2M$ it is the kinetic energy $\tfrac12\vect{p}\cdot(M + \delta M)^{-1}\cdot\vect{p}$ to first order. For an atom the adiabatic states depend on the position $\vect{X}$ of the nucleus only through $\vect{r}_e - \vect{X}$, and \cref{sec:6.1} showed that $\vect{d}_{j0} = m_e\Delta E_j\langle\Psi_j|\sum_e\vect{r}_e|\Psi_0\rangle$. Then $\vect{F}_{j0} = -m_e\Delta E_j^2\langle\Psi_j|\sum_e\vect{r}_e|\Psi_0\rangle$, and a diagonal element of \cref{eq:D.28-dM} is $2m_e^2\sum_j\Delta E_j\allowbreak|\langle\Psi_j|\sum_ex_e|\Psi_0\rangle|^2 = m_e\sum_jf^{\mathrm{osc}}_{j0}$ with the oscillator strengths of \cref{eq:4.3-osc}. By the Thomas--Reiche--Kuhn rule \cref{eq:4.3-TRK}, $\delta M = N_em_e$ for an atom with $N_e$ electrons: the whole cloud is dragged, and the atom moves with its full mass, as in the solvable model of \cref{app:D.13}.

\subsection{The kick and the electronic sectors}

Consider one nucleus at $\vect{X}$ and the kick $\ee^{\ii\vect{q}\cdot\vect{X}}$, which is diagonal in the electronic label at each $\vect{X}$. The tilted initial state is $\sum_k\Psi^{\mathrm{nuc}}_k\Psi_k$, with $\Psi^{\mathrm{nuc}}_k$ from \cref{eq:D.28-tilt}. The tilted final state on surface $j$ is $\sum_k\Psi^{\mathrm{nuc}\prime}_k\Psi_k$, dominated by its vibrational state $\Psi^{\mathrm{nuc}\prime}_j$, and its component on the ground surface is, by the same argument with $j$ and 0 exchanged, $\Psi^{\mathrm{nuc}\prime}_0 = \vect{d}_{0j}\cdot\nabla\Psi^{\mathrm{nuc}\prime}_j/[M(E_0 - E_j)] = \vect{d}_{j0}^*\cdot\nabla\Psi^{\mathrm{nuc}\prime}_j/(M\Delta E_j)$, using $\vect{d}_{0j} = -\vect{d}_{j0}^*$, the anti-Hermiticity of the derivative coupling; the complex conjugate of $\Psi^{\mathrm{nuc}\prime}_0$ that enters the overlap below then carries $\vect{d}_{j0}$ itself. The leading term $\langle\Psi^{\mathrm{nuc}\prime}_j\Psi_j|\ee^{\ii\vect{q}\cdot\vect{X}}|\Psi^{\mathrm{nuc}}_0\Psi_0\rangle$ vanishes by orthogonality of the adiabatic states at each $\vect{X}$, and to first order the amplitude is
\begin{multline*}
  \langle\Psi^{\mathrm{nuc}\prime}_j|\ee^{\ii\vect{q}\cdot\vect{X}}|\Psi^{\mathrm{nuc}}_j\rangle + \langle\Psi^{\mathrm{nuc}\prime}_0|\ee^{\ii\vect{q}\cdot\vect{X}}|\Psi^{\mathrm{nuc}}_0\rangle\\
  = \frac{1}{M\Delta E_j}\int\dd\vect{X}\,\Bigl[\Psi^{\mathrm{nuc}\prime*}_j\ee^{\ii\vect{q}\cdot\vect{X}}\,\vect{d}_{j0}\cdot\nabla\Psi^{\mathrm{nuc}}_0 + (\vect{d}_{j0}\cdot\nabla\Psi^{\mathrm{nuc}\prime*}_j)\,\ee^{\ii\vect{q}\cdot\vect{X}}\Psi^{\mathrm{nuc}}_0\Bigr] .
\end{multline*}
In the second term we integrate by parts. Treating $\vect{d}_{j0}$ and $\Delta E_j$ as constant over the extent of the nuclear wavefunctions, which neglects corrections of relative order $\ell_{\mathrm{zp}}|\nabla\ln\vect{d}_{j0}| \sim \kappa_{\mathrm{BO}}$, the gradient falls on $\ee^{\ii\vect{q}\cdot\vect{X}}\Psi^{\mathrm{nuc}}_0$ and gives $-\Psi^{\mathrm{nuc}\prime*}_j\ee^{\ii\vect{q}\cdot\vect{X}}(\ii\vect{q}\cdot\vect{d}_{j0}\Psi^{\mathrm{nuc}}_0 + \vect{d}_{j0}\cdot\nabla\Psi^{\mathrm{nuc}}_0)$. The gradients of $\Psi^{\mathrm{nuc}}_0$ cancel between the two terms, and
\begin{equation}
  \bigl\langle\widetilde{\Psi^{\mathrm{nuc}\prime}\Psi_j}\big|\ee^{\ii\vect{q}\cdot\vect{X}}\big|\widetilde{\Psi^{\mathrm{nuc}}\Psi_0}\bigr\rangle = -\ii\,\frac{\vect{q}\cdot\vect{d}_{j0}}{M\Delta E_j}\,\langle\Psi^{\mathrm{nuc}\prime}_j|\ee^{\ii\vect{q}\cdot\vect{X}}|\Psi^{\mathrm{nuc}}_0\rangle ,
  \label{eq:D.28-boost}
\end{equation}
which is \cref{eq:8.6-boost}. The cancellation makes the result independent of how the nucleus moved before the kick; only the momentum transferred appears. Summing the squared modulus over the complete set of vibrational states $\Psi^{\mathrm{nuc}\prime}_j$ of surface $j$ removes the nuclear factor, since $\sum|\langle\Psi^{\mathrm{nuc}\prime}_j|\ee^{\ii\vect{q}\cdot\vect{X}}|\Psi^{\mathrm{nuc}}_0\rangle|^2 = 1$, and gives $P_{\mathrm{Mig}}$ of \cref{eq:8.6-PMig}.

\subsection{Partial sum rules}

At $T \ll \Eg$ the initial states lie in the electron-following sector. For a projector that commutes with $H$, the response of sector ad is the full response of the restricted operators $\bar O = \mathcal{P}_{\mathrm{ad}}O\mathcal{P}_{\mathrm{ad}}$ under $H_{\mathrm{ad}} = \mathcal{P}_{\mathrm{ad}}H\mathcal{P}_{\mathrm{ad}}$, and the first-moment identity \cref{eq:4.3-firstmoment} applies to it inside the sector: $\mathscr{M}^1_{ab}|_{\mathrm{ad}} = \Vol^{-1}\langle[\bar O_a^\dagger, [H_{\mathrm{ad}}, \bar O_b]]\rangle$. Rotated by $U_n^\dagger$, the sector becomes functions of $u$ times $\Psi_0(u)$, $H_{\mathrm{ad}}$ becomes $h_{\mathrm{BO}}$, and the electron density becomes $\bar\rho_e(\vect{q};u)$ of \cref{eq:8.6-follow} up to the tilt, of relative order $\kappa_{\mathrm{BO}}^2$. A function $F(u)$ commutes with $E_0(u)$, and the double commutator with the kinetic energy of the nuclei is $[F^\dagger, [p^2/2M, F]] = |\nabla F|^2/M$, from $[p^2, F] = -\nabla^2F - 2\ii(\nabla F)\cdot\vect{p}$. Summed over nuclei,
\begin{equation*}
  \mathscr{M}^1_{ee}\big|_{\mathrm{ad}} = \frac1\Vol\Bigl\langle\sum_{\vect{R}\kappa}\frac{|\nabla_{\vect{R}\kappa}\bar\rho_e|^2}{M_\kappa}\Bigr\rangle , \qquad
  \mathscr{M}^1_{e\kappa}\big|_{\mathrm{ad}} = \frac1\Vol\Bigl\langle\sum_{I\in\kappa}\frac{\nabla_I\bar\rho_e^\dagger\cdot\nabla_I\ee^{\ii\vect{q}\cdot\vect{X}_I}}{M_\kappa}\Bigr\rangle ,
\end{equation*}
with $\nabla_I\ee^{\ii\vect{q}\cdot\vect{X}_I} = \ii\vect{q}\,\ee^{\ii\vect{q}\cdot\vect{X}_I}$. The mass correction changes these by the relative amount $\delta M/M \sim \kappa_{\mathrm{BO}}^4$. For rigid clouds, $\bar\rho_e = \sum_If_{\kappa(I)}(\vect{q})\ee^{\ii\vect{q}\cdot\vect{X}_I}$, the gradient with respect to $\vect{X}_I$ is $\ii\vect{q}f_\kappa\ee^{\ii\vect{q}\cdot\vect{X}_I}$, and the averages give $q^2\sum_\kappa\nd_\kappa|f_\kappa|^2/M_\kappa$ and $q^2\nd_\kappa f_\kappa^*/M_\kappa$, the forms quoted in \cref{eq:8.6-fsum-e,eq:8.6-fsum-ek}. The total of $\mathscr{M}^1_{e\kappa}$ over all sectors is zero by \cref{eq:4.3-species}, and the electronic sectors carry the negative of the electron-following part.

For the nuclear density the electronic share follows from \cref{eq:D.28-boost}. Each final state lies above the initial one by $\Delta E_j$ plus the vibrational energy of $\Psi^{\mathrm{nuc}\prime}_j$ relative to $\Psi^{\mathrm{nuc}}_0$. Weighted with the squared amplitude \cref{eq:D.28-boost} and summed over the vibrational states, the electronic energies give $\sum_j|\vect{q}\cdot\vect{d}_{j0}|^2/(M^2\Delta E_j) = \sum_j|\vect{q}\cdot\vect{F}_{j0}|^2/(M^2\Delta E_j^3) = \vect{q}\cdot\delta M\cdot\vect{q}/2M^2$, and the vibrational energies give the first moment of the kick on the excited surface, $E_R$ by the argument of \cref{eq:5.6-sum}, times $P_{\mathrm{Mig}}$. The first moment summed over $\pm\vect{q}$ is twice the one-sided one, and per nucleus the electronic sectors carry $\mathscr{M}^1|_{0+} = \nd_\kappa[P_{\mathrm{Mig}}q^2/M + \vect{q}\cdot\delta M\cdot\vect{q}/M^2]$. The total is $\nd_\kappa q^2/M$, and the electron-following sector carries the difference, $\nd_\kappa[(1 - P_{\mathrm{Mig}})q^2/M - \vect{q}\cdot\delta M\cdot\vect{q}/M^2]$. To first order in the two small quantities this is $\nd_\kappa(1 - P_{\mathrm{Mig}})\,\vect{q}\cdot(M + \delta M)^{-1}\cdot\vect{q}$, the same result that follows from $h_{\mathrm{BO}}$ with the depleted weight $1 - P_{\mathrm{Mig}}$ of $\bar n_\kappa$. This is \cref{eq:8.6-fsum-k}.

\subsection{Labels inside the off-shell window}

Let $H_0$ contain the electronic excitations at clamped, renormalized bands and the harmonic phonons, and let $H_{\mathrm{ep}}$ couple them; the eigenstates of $H_0$ carry a definite phonon number. Take $\omega$ inside the off-shell window, so that every zero-phonon state $|m\rangle$ at the same total momentum has $E_m \ge \omega + \delta_{\mathrm{w}}$. An exact final state at $\omega$ grows out of a one-phonon state $|f\rangle$ of $H_0$. The Lippmann--Schwinger equation $|\tilde f\rangle = |f\rangle + (E_f - H_0 + \ii0^+)^{-1}H_{\mathrm{ep}}|\tilde f\rangle$, projected on the zero-phonon states, has denominators $E_f - E_m \le -\delta_{\mathrm{w}}$ that never vanish, and its iteration converges for $|g| \ll \delta_{\mathrm{w}}$. The zero-phonon component is $\sum_m|m\rangle g_{mf}/(E_f - E_m) + O(g^3/\delta_{\mathrm{w}}^3)$, with squared norm $\sum_m|g_{mf}|^2/(E_f - E_m)^2 \le \theta_{\mathrm{ep}}$ when $\delta_{\mathrm{off}}$ in \cref{eq:8.5-theta} is taken as the smallest denominator, $\delta_{\mathrm{w}}$. Any two definitions of the label that agree at $g = 0$ and are continuous in $g$, the phonon number of the leading component, a Schrieffer--Wolff number, the number of phonons in the outgoing state, assign $|\tilde f\rangle$ to $\mathcal{P}_+$ and differ by at most this weight. The response at $\omega$ is $(2\pi/\Vol)\sum_f|\langle\tilde f|O_b|i\rangle|^2\delta(\omega - E_f)$, all of it in $\mathcal{P}_+$ to relative order $\theta_{\mathrm{ep}}$. To leading order $\langle\tilde f|O_b|i\rangle = g_{fm}[\mathcal{A}^{\mathrm{off}}_f + \rho_m/(\omega - E_m)]$, summed over $m$, which is \cref{eq:8.6-dyson} with no delta term, since $\omega - \tilde E_m$ never vanishes on the support. The bare projector $\mathcal{P}^{(0)}_0$, by contrast, gives $\langle\tilde f|\mathcal{P}^{(0)}_0O_b|i\rangle = \sum_mg^*_{mf}\rho_m/(E_f - E_m)$, the resonant part of the same amplitude, and it would assign to $\mathcal{P}_0$ a weight of the size of the rate.

\subsection{The resonant amplitude}

Outside the window, keep one zero-phonon pair $|m\rangle$ at energy $E_m$ and the one-phonon states $|f\rangle$ it couples to, with $H = E_m|m\rangle\langle m| + \sum_fE_f|f\rangle\langle f| + \sum_f(g_{fm}|f\rangle\langle m| + g_{fm}^*|m\rangle\langle f|)$; nonresonant pairs and the paths through occupied states only enter the amplitudes $\mathcal{A}^{\mathrm{off}}_f$, and degenerate pairs need the matrix form of the same algebra. The partitioning \cref{eq:C.7-partition} with $P = |m\rangle\langle m|$ gives the level's Green's function $\GF_m(z) = [z - E_m - \Sigma_m(z)]^{-1}$ with $\Sigma_m(z) = \sum_f|g_{fm}|^2/(z - E_f)$, as in \cref{we:4-level}. The scattering states $|f^-\rangle$ of the continuum satisfy $\langle f^-|m\rangle = g_{fm}\GF_m(E_f + \ii0^+)$ and $\langle f^-|f'\rangle = \delta_{ff'} + g_{fm}\GF_m(E_f + \ii0^+)\,g_{f'm}^*/(E_f - E_{f'} + \ii0^+)$, the Lippmann--Schwinger solution of the model~\cite{fano1961effects}. If no bound state splits off, they are complete, and the response of the source $O_b|i\rangle = \rho_m|m\rangle + \sum_fg_{fm}\mathcal{A}^{\mathrm{off}}_f|f\rangle$ is $(2\pi/\Vol)\sum_f|\langle f^-|O_b|i\rangle|^2\delta(\omega - E_f)$. The amplitude is \cref{eq:8.6-amp}.

Let $\mathcal{A}^{\mathrm{off}}_f$ take the value $\mathcal{A}^{\mathrm{off}}_\nu$ on the energy shell of branch $\nu$, and use $\sum_{f\in\nu}|g_{fm}|^2\delta(\omega - E_f) = \varGamma^{(\nu)}_m(\omega)/2\pi$. The response in branch $\nu$ is $(2\pi/\Vol)(\varGamma^{(\nu)}_m/2\pi)|\mathcal{A}^{\mathrm{off}}_\nu + \GF_m(\rho_m + \Upsilon_m)|^2$, and we expand its three terms near $\tilde E_m$, with $x = \omega - \tilde E_m$, against a test function smooth on $\delta_{\mathrm{F}}$. Near the pole $\GF_m \simeq z_{\mathrm{qp}}/(x + \ii z_{\mathrm{qp}}\varGamma_m/2)$, and \cref{eq:2.4-split} with the width $z_{\mathrm{qp}}\varGamma_m$ gives $|\GF_m|^2 = (2\pi z_{\mathrm{qp}}/\varGamma_m)\delta(x) + \operatorname{fp}(1/x^2) + O(g^2)$. In the cross term $2\Real[\mathcal{A}^{\mathrm{off}*}_\nu\GF_m\rho_m]$ the propagator becomes $\operatorname{PV}(1/x) - \ii\pi\delta(x)$ by \cref{eq:C.4-SP}, which yields $2\Real(\mathcal{A}^{\mathrm{off}*}_\nu\rho_m)\operatorname{PV}(1/x) + 2\pi\Imag(\mathcal{A}^{\mathrm{off}*}_\nu\rho_m)\delta(x)$. In $|\rho_m + \Upsilon_m|^2 = |\rho_m|^2 + 2\Real(\rho_m^*\Upsilon_m) + O(g^4)$, the vertex correction at the pole is $\Upsilon_m = \operatorname{PV}\Upsilon_m - (\ii/2)\sum_\mu\varGamma^{(\mu)}_m\mathcal{A}^{\mathrm{off}}_\mu$, and $2\Real(\rho_m^*\Upsilon_m) = 2\Real(\rho_m^*\operatorname{PV}\Upsilon_m) + \sum_\mu\varGamma^{(\mu)}_m\Imag(\rho_m^*\mathcal{A}^{\mathrm{off}}_\mu)$. With $c_\mu = \Imag(\rho_m^*\mathcal{A}^{\mathrm{off}}_\mu)$ and $\bar c = \sum_\mu\varGamma^{(\mu)}_mc_\mu/\varGamma_m$, the coefficient of $(2\pi/\Vol)\delta(x)$ in branch $\nu$ collects to
\begin{equation}
  \frac{\varGamma^{(\nu)}_m}{\varGamma_m}\Bigl[z_{\mathrm{qp}}|\tilde\rho_m|^2 + \varGamma_m(\bar c - c_\nu)\Bigr]
  = z_{\mathrm{qp}}|\tilde\rho_m|^2\,\frac{\varGamma^{(\nu)}_m}{\varGamma_m}\bigl(1 + \bar c^{\mathrm{F}} - c^{\mathrm{F}}_\nu\bigr) + O(g^4) ,
  \label{eq:D.28-delta}
\end{equation}
with $c^{\mathrm{F}}_\nu = \varGamma_mc_\nu/|\rho_m|^2$. Summed over $\nu$ with the weights $\varGamma^{(\nu)}_m/\varGamma_m$, which add up to one, the Fano terms cancel, and the delta part is $S^0$ of \cref{eq:8.6-dyson}. The rest, $(\varGamma^{(\nu)}_m/2\pi)[|\rho_m|^2\operatorname{fp}(1/x^2) + 2\Real(\mathcal{A}^{\mathrm{off}*}_\nu\rho_m)\operatorname{PV}(1/x) + |\mathcal{A}^{\mathrm{off}}_\nu|^2]$, summed over $\nu$, is $S^+$. The neglected terms are of three kinds: order $g^4$, from $z_{\mathrm{qp}}^2 \to 1$ in the finite part and from $|\Upsilon_m|^2$; the term $-\pi\varGamma_m\delta''(x)/4$ of \cref{eq:2.4-split}, which against a test function smooth on $\delta_{\mathrm{F}}$ is of relative order $\varGamma_m^2/\delta_{\mathrm{F}}^2$; and the shift of the finite part from $E_m$ to $\tilde E_m$, of order $g^4$. The totals therefore agree with the exact response to $O(g^4, \varGamma_m^2/\delta_{\mathrm{F}}^2)$. If the pair also decays by channels other than phonon emission, impact ionization for instance, the same algebra holds with $\nu$ and $\mu$ running over every channel. Each channel has its own off-shell amplitude $\mathcal{A}^{\mathrm{off}}_\nu$, the self-energy and the vertex correction $\Upsilon_m$ sum over all of them, the Lorentzian carries the total width $\varGamma^{\mathrm{tot}}_m = \sum_\nu\varGamma^{(\nu)}_m$, and \cref{eq:D.28-delta} becomes \cref{eq:8.6-fano}. The deviations $c^{\mathrm{F}}_\nu - \bar c^{\mathrm{F}}$ then average to zero over all channels; restricted to the phonon branches but normalized by $\varGamma^{\mathrm{tot}}_m$, they would not, and the number of zero-phonon events would change. As before, the delta function itself needs no width, and only the cascade needs $\varGamma^{\mathrm{tot}}_m$. A cascade that lets the zero-phonon pair emit its first phonon with these branching ratios, Fano terms included, reproduces the exact distribution of the first emission at order $g^2$.

The residue identity \cref{eq:8.6-zsum} follows from the Kramers--Kronig form $\Real\Sigma_m(\omega) = (2\pi)^{-1}\operatorname{PV}\!\int\dd\omega'\,\varGamma_m(\omega')/(\omega - \omega')$ and \cref{eq:C.4-dPV}. The ambiguity estimate needs a bound on $1 - z_{\mathrm{qp}}^{-1} = \partial_\omega\Real\Sigma_m$. Split the finite-part integral at $|\omega' - \tilde E_m| = \delta_{\mathrm{F}}$. Outside, $|\int\dd\omega'\,\varGamma_m/(\omega' - \tilde E_m)^2| \le 2\hat\varGamma_m/\delta_{\mathrm{F}}$, with $\hat\varGamma_m$ the largest width. Inside, $\varGamma_m(\tilde E_m)$ times $\operatorname{fp}\!\int_{-\delta_{\mathrm{F}}}^{\delta_{\mathrm{F}}}\dd y/y^2 = -2/\delta_{\mathrm{F}}$ contributes $2\varGamma_m/\delta_{\mathrm{F}}$ in modulus, the linear term of $\varGamma_m$ drops out by symmetry, and the curvature contributes at most $\delta_{\mathrm{F}}\|\varGamma_m''\|$, which is $\hat\varGamma_m/\delta_{\mathrm{F}}$ if $\varGamma_m$ varies on the scale $\delta_{\mathrm{F}}$. Hence $|1 - z_{\mathrm{qp}}^{-1}| \le (5/2\pi)\hat\varGamma_m/\delta_{\mathrm{F}}$. Raising the residue to a power between 1 and 2 changes the zero-phonon weight by at most this amount, relative to $w_0$, which is the first term of \cref{eq:8.6-ambiguity}. Replacing the delta function by a normalized kernel of width $\varGamma_m$ moves the tail weight $2\int_{x_{\mathrm{w}}}^\infty\dd x\,(\varGamma_m/2\pi)/x^2 = \varGamma_m/\pi x_{\mathrm{w}}$ outside a window of half-width $x_{\mathrm{w}}$, the second term. Assigning to the line all weight within $|x| < x_{\mathrm{w}}$ moves into it the finite-part, principal-value and $|\mathcal{A}^{\mathrm{off}}|^2$ terms of $S^+$ there. Their pole parts reproduce the second term; their smooth parts are bounded by $\hat\varGamma_mx_{\mathrm{w}}/\delta_{\mathrm{F}}^2$ times $|\rho_m|^2$ for the finite part, $2|\rho_m||\mathcal{A}^{\mathrm{off}}|\delta_{\mathrm{F}}$ for the principal value and $|\mathcal{A}^{\mathrm{off}}|^2\delta_{\mathrm{F}}^2$ for the last term. Divided by $w_0 \simeq |\rho_m|^2$ they add up to the third term. Each step is first order in $\varGamma_m/\delta_{\mathrm{F}}$, and the bound holds to that order.

\subsection{Real amplitudes under inversion and time reversal}

Let $\mathsf{P}\mathsf{T}$ be the product of inversion and time reversal. It is antiunitary and commutes with $H$ if both symmetries hold. On states of $N$ electrons $\mathsf{T}^2 = (-1)^N$, $\mathsf{P}^2 = 1$ and the two commute, so $(\mathsf{P}\mathsf{T})^2 = (-1)^N$; on phonons $(\mathsf{P}\mathsf{T})^2 = 1$. For even $N$, then, $(\mathsf{P}\mathsf{T})^2 = 1$ on the whole space. An antiunitary involution that commutes with a Hermitian operator leaves each eigenspace invariant, and every eigenspace has an orthonormal basis of invariant vectors: for any $v$ in it, $v + \mathsf{P}\mathsf{T}v$ and $\ii(v - \mathsf{P}\mathsf{T}v)$ are invariant, and $v$ is a combination of them. Choose such a basis for the eigenstates of $H_0$ within each sector, which is possible because $\mathsf{P}\mathsf{T}$ commutes with the phonon number. For the one-phonon states this fixes the phase of each polarization vector $\vect{e}_{\kappa\nu}(\vect{Q})$, since inversion and time reversal each reverse $\vect{Q}$ and $\mathsf{P}\mathsf{T}$ maps the mode $(\vect{Q}, \nu)$ to itself. For two invariant states $a$, $b$ and an invariant operator $O$, antiunitarity gives $\langle a|O|b\rangle = \langle\mathsf{P}\mathsf{T}a|\mathsf{P}\mathsf{T}Ob\rangle^* = \langle a|Ob\rangle^*$, a real number. The density is invariant: inversion maps $n_e(\vect{q})$ to $n_e(-\vect{q})$, and time reversal, which conjugates $\ee^{-\ii\vect{q}\cdot\vect{r}}$, maps it back. So is $H_{\mathrm{ep}}$, as part of $H$. Every $\rho_m$ and $g_{fm}$ is real, and so is every off-shell amplitude $\mathcal{A}^{\mathrm{off}}_f$, a sum of products of real matrix elements and real denominators. Hence $c_\nu = 0$ and the Fano coefficients vanish, provided the invariant basis of a degenerate zero-phonon manifold also diagonalizes its level-shift matrix, which holds for one-body self-energies and can fail for level shifts from the electron--hole vertex (\cref{sec:8.6}). With spin--orbit coupling the single-electron states form Kramers doublets with $(\mathsf{P}\mathsf{T})^2 = -1$, but the argument concerns many-electron states with even $N$ and needs no real single-particle gauge.

\subsection{Gapless targets}

For a state $\psi_b$ that decays by the golden rule into a continuum of final states $f$ with a flat density $\mathcal{D}_f$ over a width $\Delta E_{\mathrm{c}}$, $\varGamma = 2\pi|\langle f|H_{\mathrm{off}}|\psi_b\rangle|^2\mathcal{D}_f/\|\psi_b\|^2$ and $\|H_{\mathrm{off}}\psi_b\|^2 = \sum_f|\langle f|H_{\mathrm{off}}|\psi_b\rangle|^2 = \varGamma\,\Delta E_{\mathrm{c}}\|\psi_b\|^2/2\pi$, which is the estimate of \cref{sec:8.6}. Without any dressing, $U_{\mathrm{qp}} = 1$, the estimate does not become a bound. The diagonal part $H_{\mathrm{d}}$ conserves the bare numbers, but the probe's quasiparticle is created on top of the filled Fermi sea, and $H_{\mathrm{off}}$ also acts on the sea, creating a phonon and a pair away from the probe's excitations. The squared norm of this disconnected term grows with the number of cells, and so does $\mathfrak{h}_{\mathrm{off}}$. The estimate describes the connected part, in which $H_{\mathrm{off}}$ acts on the quasiparticle $|n\vect{k}\rangle$ itself. For the electron--phonon part, the state $H_{\mathrm{ep}}|n\vect{k}\rangle$ contains one phonon emitted or absorbed, and its squared norm is $N^{-1}\sum_{m\nu\vect{Q}}|g_{mn\nu}(\vect{k},\vect{Q})|^2$ times $n_{\vect{Q}\nu} + 1 - f_{m,\vect{k}+\vect{Q}}$ for emission and $n_{\vect{Q}\nu} + f_{m,\vect{k}+\vect{Q}}$ for absorption, at most $2n_{\vect{Q}\nu} + 1$ in total. This is $-\pi^{-1}\int\dd\omega\,\Imag\Sigma^{\mathrm{FM}}_{n\vect{k}}(\omega)$ from \cref{eq:6.3-ImFM}, the zeroth moment of the Fan--Migdal spectral function. The electron--electron part adds the analogous sum over final states with two particles and a hole and the phonon number unchanged; the two sets of final states differ in phonon number, so they are orthogonal and the sums add. For a pair source, the electron and hole contributions add. Taken over the bands, these moments form a matrix whose largest eigenvalue bounds the square of the connected part. In aluminum the electron--phonon moments alone give an eigenvalue of about \SI{0.03}{eV^2} (\cref{sec:8.6}), so the electron--phonon contribution to the connected part is at most \SI{0.17}{eV}; the electron--electron moments cannot lower the eigenvalue and are not estimated. Separating the connected part from the disconnected one needs the dressed vacuum $U_{\mathrm{qp}}$ and a commutator bound uniform in $N$, which is open (\cref{sec:8.6}).
